\section{Signed projected forcing and the original upper endpoint}
\label{sec:v94-signed}
Taking positive parts on every sheet can overestimate the forcing of
the projected source. We retain all physical and prefix losses through
projection. Only then is a positive part taken. The resulting forcing
is smaller than the earlier sufficient positive budgets and gives an
equivalent criterion for the original scalar endpoint.

Restrict $\mathsf G_c^{\varepsilon,+}$ at its physical intermediate
state to $z<\rho$ or $z\ge\rho$ and set
\begin{equation}\label{eq:v94-net-measure}
 \mathsf K_{c,\rho}^\varepsilon
 =\mathsf I_c^\varepsilon+\mathsf G_c^{\varepsilon,+,<\rho}
               -\mathsf D_c^\varepsilon-\mathsf G_c^{\varepsilon,-}.
\end{equation}
This is a finite signed measure at fixed count. For such a measure put
\[
 \mathcal T_r\mu(t)=\int_{(t-r,t]}k_{r,t}(u)\,d\mu(u),\qquad
 A_r^-f(t)=\frac1r\int_{t-r}^t f(u)\,du.
\]
Define two nonnegative functions
\begin{align}
 S_{c;r,\rho}^\varepsilon(t)
    &=[\mathcal T_r\mathsf K_{c,\rho}^\varepsilon(t)]_+,\notag\\
 S_{c;r,\rho}^{\varepsilon,\sharp}(t)
    &=\operatorname*{ess\,sup}_{s\in(t-r,t)}
              [\mathsf K_{c,\rho}^\varepsilon((s,t])]_+.
                                             \label{eq:v94-signed-forcing}
\end{align}
The supremum is in the earlier roof only. BV representatives and
Fubini give Borel versions of both functions; pointwise assertions
below are at Lebesgue roofs. They are not Jordan decompositions of
$D R_c^\varepsilon$.

\begin{theorem}[A complete signed pointwise comparison]
\label{thm:v94-signed-comparison}
For every fixed $r,\rho>0$ and almost every roof,
\begin{align}
 0\le S_{c;r,\rho}^\varepsilon
       &\le S_{c;r,\rho}^{\varepsilon,\sharp}\le R_c^\varepsilon,
                                            \label{eq:v94-forcing-domination}\\
 R_c^\varepsilon(t)
       &\le A_r^-R_c^\varepsilon(t)+S_{c;r,\rho}^\varepsilon(t)
         +C(1+\varepsilon^{-1})\int_{t-r}^t v^\rho(u)\,du.
                                            \label{eq:v94-signed-pointwise}
\end{align}
For almost every $s\in(t-r,t)$ one also has
\begin{equation}\label{eq:v94-signed-persistence}
 R_c^\varepsilon(s)\ge R_c^\varepsilon(t)
       -S_{c;r,\rho}^{\varepsilon,\sharp}(t)
       -C(1+\varepsilon^{-1})\int_{t-r}^t v^\rho(u)\,du.
\end{equation}
Completed prefixes contribute nonpositively before the positive part
in \eqref{eq:v94-signed-forcing}. Moreover
\begin{equation}\label{eq:v94-smaller-forcing}
 S_{c;r,\rho}^\varepsilon(t)
 \le B_{c,r}^\varepsilon(t)+G_{c,r}^{\varepsilon,+,<\rho}(t),
\end{equation}
where the right side is the retained target-ancestral ramped budget.
The analogous inequality holds with all ramps removed.
\end{theorem}
\begin{proof}
By the exact signed balance,
$D R_c^\varepsilon=\mathsf K_{c,\rho}^\varepsilon+
\mathsf G_c^{\varepsilon,+,\ge\rho}$. Stieltjes integration gives
\begin{equation}\label{eq:v94-exact-ramp}
 R_c^\varepsilon(t)-A_r^-R_c^\varepsilon(t)
   =\mathcal T_r\mathsf K_{c,\rho}^\varepsilon(t)
       +\mathcal T_r\mathsf G_c^{\varepsilon,+,\ge\rho}(t).
\end{equation}
The last term is nonnegative. Since $A_r^-R_c^\varepsilon\ge0$,
its subtraction proves $S_{c;r,\rho}^\varepsilon\le R_c^\varepsilon$.
For almost every earlier roof, the interval form of the same identity is
\[
 \mathsf K_{c,\rho}^\varepsilon((s,t])
 =R_c^\varepsilon(t)-R_c^\varepsilon(s)
       -\mathsf G_c^{\varepsilon,+,\ge\rho}((s,t])
 \le R_c^\varepsilon(t).
\]
It proves the sharp domination and the persistence bound. Averaging
this interval identity in $s$ gives \eqref{eq:v94-exact-ramp}; hence
the ramped positive part is bounded by the sharp one. Finally, the
full physical product derivative and multiplicity-one source sum give
\[
 \mathsf G_c^{\varepsilon,+,\ge\rho}((t-r,t])
       \le C(1+\varepsilon^{-1})\int_{t-r}^t v^\rho(u)\,du.
\]
This proves \eqref{eq:v94-signed-pointwise} and
\eqref{eq:v94-signed-persistence}.

For the finer comparison, do the same calculation on one zero-extended
prefix. If $c_\sigma\le t$, its ramped derivative is
$-A_r^-(\one_{(a_\sigma,c_\sigma)}f_\sigma^\varepsilon)(t)\le0$;
subtracting its noncritical positive growth only decreases it.
A future prefix contributes zero. For a live target ancestor there
is no upper cut in the window, and its remaining signed contribution
is at most its entrance trace plus critical positive growth. Sum all
signed contributions before taking the positive part. Fixed-count
finite variation justifies the interchange. This proves
\eqref{eq:v94-smaller-forcing}; the interval form proves the sharp
version. In particular no unobserved completed history is charged as
positive forcing.
\end{proof}

\subsection{Where the signed forcing can still be positive}
Write $R_c^{\varepsilon,<2b}$ for the original prefix source restricted
at the target to $z(t)<2b$, and let $B_{c,r}^{\varepsilon,b,\sharp}$
be its entrance budget restricted to $z_\sigma^-\ge b$, without a
residence cutoff.

\begin{proposition}[Critical-target and inlet localization]
\label{prop:v94-localization}
For $0<\rho\le b$ and the same $r,b$ as in
Theorem~\ref{thm:v94-transit},
\begin{equation}\label{eq:v94-localization}
 S_{c;r,\rho}^{\varepsilon,\sharp}(t)
 \le R_c^{\varepsilon,<2b}(t)
              +B_{c,r}^{\varepsilon,b,\sharp}(t)
              +N_{c;r,b}^\varepsilon(t).
\end{equation}
Thus radial-transiting critical packets are already estimated, and
all noncritical inlet residence bands above a fixed positive cutoff
are already estimated. The two remaining terms are restrictions of
the original critical target source and original zero-residence inlet
source, not separate absolute critical-growth charges.
\end{proposition}
\begin{proof}
Completed or future prefixes contribute nonpositively to every earlier
interval comparison. For live prefixes, first separate targets with
$z(t)<2b$. Each sheet's signed interval increment is at most its
current value, so their total positive contribution is at most
$R_c^{\varepsilon,<2b}(t)$. For the remaining targets use
$u_\sigma^0=\max(a_\sigma,t-r)$. If $z(u_\sigma^0+)<b$, the sheet
belongs to the radial-transit family and its signed increment is at
most $f_\sigma^\varepsilon(t)$. If $z(u_\sigma^0+)\ge b$, no critical
positive growth occurs in the window. An old sheet then contributes
nonpositively; a newborn sheet contributes at most its entrance trace.
These bounds hold for every earlier roof outside the common null set,
so they survive the earlier-roof essential supremum. All source
restrictions are Borel and all sheet interiors are counted once.
\end{proof}

\subsection{An equivalent upper-endpoint criterion}
For the full source $c_\sigma=b_\sigma$ omit $c$ and define
\[
 \mathfrak S_{\varepsilon,M}(r)
  =\lim_{\rho\downarrow0}\mathcal H_M(S_{r,\rho}^\varepsilon),
 \qquad
 \mathfrak S_{\varepsilon,M}^{\sharp}(r)
  =\lim_{\rho\downarrow0}\mathcal H_M(S_{r,\rho}^{\varepsilon,\sharp}).
\]
Both limits exist monotonically: decreasing $\rho$ removes a positive
part from \eqref{eq:v94-net-measure}. The complete finite-count signed
sum, its positive part and its essential height all precede this limit.

\begin{theorem}[Signed forcing at the original pointwise topology]
\label{thm:v94-endpoint}
For all large fixed $B$ and fixed $r>0$,
\begin{equation}\label{eq:v94-endpoint}
 \mathcal E_M\le C_MB^{-1/192}
       +\mathfrak S_{\varepsilon(B),M}(r)
       +C(1+\varepsilon(B)^{-1})r.
\end{equation}
For any prescribed positive $r(B)$ with
$(1+\varepsilon(B)^{-1})r(B)\to0$, the original pointwise assertion
is equivalent to each of
\begin{equation}\label{eq:v94-equivalence}
 \mathfrak S_{\varepsilon(B),M}(r(B))\longrightarrow0,
 \qquad
 \mathfrak S_{\varepsilon(B),M}^{\sharp}(r(B))\longrightarrow0
                 \quad\text{for every fixed }M.
\end{equation}
The forcing in this criterion is formed after signed projection; it
is not the sum of the earlier two sufficient positive charges.
No vanishing assertion in \eqref{eq:v94-equivalence} is made here.
\end{theorem}
\begin{proof}
Apply \eqref{eq:v94-signed-pointwise} to $b^\varepsilon$.
Its backward average is at most $2A_rb^\varepsilon$. Fix
$B,r,\rho$ and all auxiliary bands before the collision limsup.
Inputs~\ref{input:v94-volume} and \ref{input:v94-scalar} then give
\[
 \mathcal H_M(b^\varepsilon)
 \le C\varepsilon^{1/16}+
      \mathcal H_M(S_{r,\rho}^\varepsilon)
                        +C(1+\varepsilon^{-1})r.
\]
Only now decrease $\rho$. Insert the original positive remainder and
$\varepsilon(B)^{1/16}=O(B^{-1/192})$ to obtain
\eqref{eq:v94-endpoint}. This proves sufficiency of either criterion,
since the ramped forcing is bounded by the sharp one.

Conversely, if the original pointwise assertion holds, the positive
remainder gives
$\mathcal H_M(b^{\varepsilon(B)})\le C_MB^{-1/192}$.
Both signed forcings are bounded by $b^{\varepsilon(B)}$ at every
finite count by \eqref{eq:v94-forcing-domination}. Their ordered
heights therefore satisfy the same bound for every $r,\rho$.
This proves necessity without exchanging any cutoff and collision
limit. The legal example $r(B)=c\varepsilon(B)^2$ makes the explicit
last term vanish, but does not prove either signed charge zero.
\end{proof}
